\begin{samepage}
    % No \chapter* here to avoid page breaks
    
    % Reduce spacing and adjust font sizes
    \begin{spacing}{1.3}
    \vspace{-1.5cm}  % Adjust vertical space if needed
    
    % Résumé section with proper formatting
    \section*{\Huge Résumé}
    \addcontentsline{toc}{section}{Résumé}  % TOC entry for Résumé
    
    \noindent
    \parbox{\textwidth}{
    Ce projet de fin d'études, réalisé au sein de SIGA, présente le développement et le déploiement d'un « Portail de Gestion des Approbations avec Workflows Dynamiques ». Face à la complexité des processus d'approbation dans les organisations, ce portail vise à centraliser et optimiser la soumission, le suivi et la validation des demandes grâce à une interface intuitive basée sur Angular, un backend robuste avec Spring Boot, et des workflows dynamiques orchestrés par Camunda. Une approche DevOps a été adoptée, intégrant un pipeline CI/CD avec GitHub Actions, la conteneurisation via Docker, et l'orchestration avec Kubernetes et ArgoCD, complétée par une gestion sécurisée des secrets. Réalisé entre janvier 2025 et mai 2025, ce projet a permis de livrer une solution scalable et fiable. Les perspectives incluent une extension à des environnements multi-clusters.
    }

    \vspace{1.4cm}  % Space between Résumé and Abstract
    
    % Abstract section with proper formatting
    \section*{\Huge Abstract}
    \addcontentsline{toc}{section}{Abstract}  % TOC entry for Abstract
    
    \noindent
    \parbox{\textwidth}{
    This final-year project, conducted at SIGA, details the development and deployment of an "Approval Management Portal with Dynamic Workflows." Addressing the complexity of approval processes in organizations, the portal aims to centralize and streamline request submission, tracking, and validation through an intuitive Angular-based interface, a robust Spring Boot backend, and dynamic workflows managed by Camunda. A DevOps approach was implemented, incorporating a CI/CD pipeline with GitHub Actions, containerization with Docker, and orchestration using Kubernetes and ArgoCD, alongside secure secret management. Developed between January 2025 and May 2025, the project delivered a scalable and reliable solution. Future prospects include expanding to multi-cluster environments.
    }
    \end{spacing}
    \end{samepage}